% ============================================================
% Part 3: Function Approximation, The Neuron & Depth,
%         Sequence Modeling, The Transformer, Scaling to LLMs
% Sections 6--10 (Slides 49--78)
% Session 3 -- The Payoff: Everything Converges
% ============================================================


% ============================================================
% SECTION 6: FUNCTION APPROXIMATION (Slides 49--53)
% "Why Neural Networks Work at All"
% Timeline: 1807--1991
% ============================================================
\section{Function Approximation}

% ----------------------------------------------------------
% Slide 49: Section Divider
% ----------------------------------------------------------
{
\setbeamercolor{background canvas}{bg=aiViolet}
\begin{frame}[plain]
  \centering
  \vspace{1.8cm}
  \textcolor{white}{\Huge\bfseries Function Approximation}\\[0.5em]
  \textcolor{white}{\Large Why Neural Networks Work at All}\\[1em]
  \textcolor{white!80}{\large 1807--1991}\\[1.5em]
  \textcolor{white!60}{\normalsize Fourier \(\cdot\) Weierstrass \(\cdot\) Cybenko \(\cdot\) Hornik}\\[1em]
  \textcolor{white!50}{\small A neural network can approximate \emph{any} function.\\
  Here is why that is remarkable.}

  \note{
    Session 3 opens here. Recap: ``We have the math (Sections 1--4) and the
    hardware (Section 5). But here is a deep question: WHY should a pile of
    matrix multiplications be able to learn ANYTHING? For the answer, we need
    the theory of function approximation.''

    Convergence diagram: Section 6 node pulsing.
  }
\end{frame}
}

% ----------------------------------------------------------
% Slide 50: Fourier Series
% ----------------------------------------------------------
\begin{frame}{Joseph Fourier --- Any Function as a Sum of Waves}

  \begin{columns}[T]
    \begin{column}{0.52\textwidth}
      \begin{personbox}[Jean-Baptiste Joseph Fourier (1768--1830)]{aiViolet}
        \textbf{Origin:} Auxerre, France\\[3pt]
        \textbf{Key work:} \emph{Th\'eorie analytique de la chaleur}
        (1822, based on work from 1807). To solve the heat equation,
        Fourier claimed that \alert{any function} could be decomposed
        into sines and cosines.\\[3pt]
        Lagrange objected. The claim was initially controversial ---
        and proved enormously fertile.
      \end{personbox}

      \vspace{0.4em}
      \begin{insightbox}
        \textbf{LLM preview:} Positional encoding in the transformer
        uses sine and cosine functions at different frequencies ---
        directly inspired by Fourier's idea. We will see this in
        Section~9.
      \end{insightbox}
    \end{column}

    \begin{column}{0.45\textwidth}
      \visualplaceholder[5.5cm]{Fourier approximation of a square wave:
      1~sine (rough), 3~sines (better), 10~sines (close),
      100~sines (nearly perfect). Caption: ``Any function = a sum
      of simple waves. Fourier, 1807.''}
    \end{column}
  \end{columns}

  \note{
    ``Fourier wanted to solve the heat equation. To do it, he made a claim
    so bold that other mathematicians did not believe him: any function can
    be written as a sum of sines. He was essentially right, and the idea
    became one of the most powerful tools in all of mathematics.''

    This slide previews positional encoding (Slide 74). Plant the seed now.
    \verifytag{Fourier's 1807 memoir date vs.\ 1822 publication date.}
  }
\end{frame}

% ----------------------------------------------------------
% Slide 51: Weierstrass Approximation [EXTENDED]
% ----------------------------------------------------------
\begin{frame}{Karl Weierstrass --- Polynomials Can Do It Too \hfill\small\textcolor{gray}{[EXTENDED]}}

  \begin{columns}[T]
    \begin{column}{0.52\textwidth}
      \begin{personbox}[Karl Weierstrass (1815--1897)]{aiViolet}
        \textbf{Origin:} Ostenfelde, Westphalia, Germany\\[3pt]
        \textbf{Key result (1885):} For any continuous function $f$ on
        a closed interval and any desired accuracy $\varepsilon > 0$,
        there exists a polynomial $P$ that stays within $\varepsilon$
        of $f$ everywhere on that interval.
      \end{personbox}

      \vspace{0.4em}
      \begin{insightbox}
        Polynomials are \textbf{universal approximators} for continuous
        functions. A century later, Cybenko would prove the same thing
        for \emph{neural networks}.
      \end{insightbox}
    \end{column}

    \begin{column}{0.45\textwidth}
      \visualplaceholder[5.5cm]{A wiggly continuous function with polynomial
      approximations of increasing degree overlaid. The polynomials get
      closer and closer. Caption: ``Any continuous function can be
      approximated as closely as you want by a polynomial.''}
    \end{column}
  \end{columns}

  \note{
    ``Weierstrass proved that polynomials are universal approximators for
    continuous functions. A century later, Cybenko would prove the same
    thing for neural networks.''

    This slide is marked EXTENDED and can be cut if short on time.
  }
\end{frame}

% ----------------------------------------------------------
% Slide 52: Universal Approximation Theorem
% ----------------------------------------------------------
\begin{frame}{The Universal Approximation Theorem}

  \begin{columns}[T]
    \begin{column}{0.55\textwidth}
      \begin{personbox}[{George Cybenko (b.\,c.\,1952)}]{aiViolet}
        \textbf{Origin:} Canada/USA\\[3pt]
        \textbf{Key result (1989):} A feedforward network with a single
        hidden layer of sigmoidal neurons can approximate \emph{any}
        continuous function to any desired accuracy.\\[3pt]
        \emph{``Approximation by Superpositions of a Sigmoidal Function''}
      \end{personbox}

      \vspace{0.3em}

      \begin{personbox}[{Kurt Hornik (b.\,c.\,1963)}]{aiViolet}
        \textbf{Origin:} Austria\\[3pt]
        \textbf{Key result (1991):} Generalized Cybenko's theorem to a
        wider class of activation functions. The universality is not
        about sigmoid --- it is about the \emph{architecture}.
      \end{personbox}
    \end{column}

    \begin{column}{0.42\textwidth}
      \visualplaceholder[5cm]{A simple neural network with one hidden layer.
      As hidden neurons increase (5, 20, 100, 1000), the function it
      represents gets more complex and accurate.
      Caption: ``More neurons = better approximation. Cybenko, 1989.''}

      \vspace{0.3em}
      \begin{tcolorbox}[colback=aiViolet!8,colframe=aiViolet!60,
        boxrule=0.5pt,rounded corners,left=3pt,right=3pt,top=2pt,bottom=2pt]
        \small A neural network \textbf{can} learn any continuous function,
        given enough neurons. This is not an analogy. It is a
        \emph{mathematical theorem}.
      \end{tcolorbox}
    \end{column}
  \end{columns}

  \note{
    ``This theorem is why people bet on neural networks. It is not that
    they are the BEST approximators --- it is that they are GUARANTEED
    to work, at least in principle. The practical challenge is training
    them, which is where all the other building blocks come in.''

    Be clear: the theorem says a network CAN approximate any function,
    not that gradient descent will FIND the right weights. The gap between
    existence and constructive algorithm is important.

    \verifytag{Cybenko birth year c.\,1952. Hornik birth year c.\,1963.}
  }
\end{frame}

% ----------------------------------------------------------
% Slide 53: Convergence Diagram Update -- Function Approximation
% ----------------------------------------------------------
\begin{frame}{Convergence Diagram --- Function Approximation Lights Up}

  \visualplaceholder[10cm]{Convergence diagram: ``Function Approximation''
  node now lit in purple. 6 of 10 nodes active.
  Label: ``Universal approximation = why it works.''}

  \vspace{0.3em}
  \begin{columns}[T]
    \begin{column}{0.48\textwidth}
      \begin{tcolorbox}[colback=aiViolet!5,colframe=aiViolet,
        fonttitle=\bfseries\small,title={Section 6 Summary},
        rounded corners,boxrule=0.8pt,left=3pt,right=3pt,top=3pt,bottom=3pt]
        \scriptsize
        \begin{tabular}{@{}ll@{}}
          Fourier (1807/1822) & Any function as waves\\
          Weierstrass (1885)  & Polynomial approximation\\
          Cybenko (1989)      & Neural net universality\\
          Hornik (1991)       & Generalized activation
        \end{tabular}
      \end{tcolorbox}
    \end{column}

    \begin{column}{0.48\textwidth}
      \begin{tcolorbox}[colback=accentOrange!8,colframe=accentOrange,
        fonttitle=\bfseries\small,
        title={\faArrowRight\ Connection to the Transformer},
        rounded corners,boxrule=0.8pt,left=3pt,right=3pt,top=3pt,bottom=3pt]
        \scriptsize
        The universal approximation theorem guarantees that neural networks
        can represent any continuous function. The LLM's ability to model
        language is a specific case.
      \end{tcolorbox}
    \end{column}
  \end{columns}

  \note{
    Transition: ``We know neural networks CAN learn anything. But what IS
    a `neural network'? Where did the idea come from?''
  }
\end{frame}


% ============================================================
% SECTION 7: THE NEURON & DEPTH (Slides 54--62)
% "From Biology to Mathematics to Deep Networks"
% Timeline: 1943--2016
% ============================================================
\section{The Neuron \& Depth}

% ----------------------------------------------------------
% Slide 54: Section Divider
% ----------------------------------------------------------
{
\setbeamercolor{background canvas}{bg=aiViolet}
\begin{frame}[plain]
  \centering
  \vspace{1.8cm}
  \textcolor{white}{\Huge\bfseries The Neuron \& Depth}\\[0.5em]
  \textcolor{white}{\Large From Biology to Deep Networks}\\[1em]
  \textcolor{white!80}{\large 1943--2016}\\[1.5em]
  \textcolor{white!60}{\normalsize McCulloch \& Pitts \(\cdot\) Rosenblatt
    \(\cdot\) Minsky \(\cdot\) Hinton \(\cdot\) He}\\[1em]
  \textcolor{white!50}{\small A real neuron fires or does not.
  A mathematical neuron does the same thing,\\with an equation.
  But stacking them deep required more invention.}

  \note{
    This section contains the most dramatic narrative arc: Perceptron hype
    $\to$ XOR crash $\to$ AI winter $\to$ backprop comeback $\to$ deep
    learning breakthroughs. Build the drama.

    Convergence diagram: Section 7 node pulsing.
  }
\end{frame}
}

% ----------------------------------------------------------
% Slide 55: The McCulloch-Pitts Neuron
% ----------------------------------------------------------
\begin{frame}{McCulloch \& Pitts --- The First Mathematical Brain Cell}

  \begin{columns}[T]
    \begin{column}{0.52\textwidth}
      \begin{personbox}[Warren McCulloch (1898--1969)]{aiViolet}
        \textbf{Origin:} Orange, New Jersey, USA\\[3pt]
        Neurophysiologist who sought to understand the brain
        in logical terms.
      \end{personbox}

      \vspace{0.3em}

      \begin{personbox}[Walter Pitts (1923--1969)]{aiViolet}
        \textbf{Origin:} Detroit, Michigan, USA\\[3pt]
        Self-taught logician. A troubled prodigy who ran away from
        home and ended up working with McCulloch.\\[3pt]
        Together they published \emph{``A Logical Calculus of the Ideas
        Immanent in Nervous Activity''} (1943).
      \end{personbox}
    \end{column}

    \begin{column}{0.45\textwidth}
      \visualplaceholder[5.5cm]{The McCulloch-Pitts neuron: inputs $x_1,
      x_2, x_3$ with weights $w_1, w_2, w_3$. Sum $=
      w_1 x_1 + w_2 x_2 + w_3 x_3$. If Sum $>$ threshold,
      output $= 1$, else output $= 0$. Biological neuron shown alongside
      for comparison.}
    \end{column}
  \end{columns}

  \note{
    ``A neurophysiologist and a teenage runaway created the mathematical
    neuron. McCulloch was in his 40s; Pitts was about 20 and had no degree.
    Their paper launched neural network research.''

    \verifytag{Pitts ``ran away at 15'' --- commonly stated but details vary.}
  }
\end{frame}

% ----------------------------------------------------------
% Slide 56: The Perceptron -- HYPE
% ----------------------------------------------------------
\begin{frame}{Frank Rosenblatt --- The Perceptron: A Machine That Learns}

  \begin{columns}[T]
    \begin{column}{0.52\textwidth}
      \begin{personbox}[Frank Rosenblatt (1928--1971)]{aiViolet}
        \textbf{Origin:} New Rochelle, New York, USA\\[3pt]
        Built the \textbf{Mark~I Perceptron} (1957--1958) at the
        Cornell Aeronautical Laboratory. A physical machine (not a
        simulation) that could learn to classify images.\\[3pt]
        The \emph{New York Times} reported it as a machine that could
        ``walk, talk, see, write, reproduce itself, and be conscious
        of its existence.''
      \end{personbox}

      \vspace{0.3em}
      \begin{insightbox}
        The perceptron learning algorithm is gradient descent
        on a single neuron. Rosenblatt's machine could \emph{learn}.
        The hype was enormous.
      \end{insightbox}
    \end{column}

    \begin{column}{0.45\textwidth}
      \visualplaceholder[5.5cm]{Photo of the Mark I Perceptron hardware:
      a large machine with photocells. Below: the perceptron learning
      algorithm --- present an example, compare output to desired output,
      adjust weights.}
    \end{column}
  \end{columns}

  \note{
    ``The press said it would think like a human. Then came the crash.''

    Build anticipation here. The next slide is the fall.

    \verifytag{Exact NYT quote and date --- commonly cited as a 1958 article.}
  }
\end{frame}

% ----------------------------------------------------------
% Slide 57: The XOR Problem and AI Winter
% ----------------------------------------------------------
\begin{frame}{The XOR Problem --- The Crash That Froze AI}

  \begin{columns}[T]
    \begin{column}{0.52\textwidth}
      \begin{personbox}[Minsky (1927--2016) \& Papert (1928--2016)]{aiViolet}
        \textbf{Marvin Minsky:} New York City, USA\\
        \textbf{Seymour Papert:} Pretoria, South Africa / USA\\[3pt]
        Published \emph{Perceptrons} (1969): a rigorous mathematical
        analysis proving that a single-layer perceptron
        \alert{cannot learn XOR}.\\[3pt]
        The math was correct. But the field interpreted it as:
        neural networks are a dead end.
      \end{personbox}

      \vspace{0.3em}
      \begin{tcolorbox}[colback=red!5,colframe=red!70,
        fonttitle=\bfseries\small,title={The AI Winter (\raisebox{0pt}{\small$\sim$}1969--1986)},
        rounded corners,boxrule=0.8pt,left=3pt,right=3pt,top=3pt,bottom=3pt]
        \small Funding dried up. Neural network research became
        unfashionable. Nearly two decades of stagnation.
        \emph{The solution was right there: add more layers.}
      \end{tcolorbox}
    \end{column}

    \begin{column}{0.45\textwidth}
      \visualplaceholder[5.5cm]{Top: XOR truth table and scatter plot showing
      XOR is not linearly separable. ``A single line cannot separate the 1s
      from the 0s. A perceptron can only draw a single line.''
      Bottom: Timeline of the AI Winter, 1969--1986.}
    \end{column}
  \end{columns}

  \note{
    ``Minsky and Papert's math was correct. But the conclusion the field
    drew --- that neural networks are a dead end --- was WRONG.
    The solution was right there: add more layers.''

    This is the dramatic low point. Pause before the next slide (the comeback).
  }
\end{frame}

% ----------------------------------------------------------
% Slide 58: Multi-Layer Networks -- The Comeback
% ----------------------------------------------------------
\begin{frame}{Multi-Layer Networks --- The Comeback}

  \begin{columns}[T]
    \begin{column}{0.55\textwidth}

      % The key point
      \begin{tcolorbox}[colback=green!5,colframe=green!60!black,
        fonttitle=\bfseries,title={1986: The Revival},
        rounded corners,boxrule=1pt,left=4pt,right=4pt,top=4pt,bottom=4pt]
        \textbf{Rumelhart, Hinton \& Williams} (Slide~22) demonstrated
        that \alert{backpropagation} could train multi-layer networks.

        \medskip
        Multi-layer perceptrons CAN learn XOR --- and much more.
        A hidden layer creates a \emph{nonlinear} decision boundary.
      \end{tcolorbox}

      \vspace{0.3em}
      \begin{insightbox}
        It took \textbf{17~years} for the field to recover from
        \emph{Perceptrons}. The AI winter was caused by stopping
        \textbf{one layer too soon}.
      \end{insightbox}
    \end{column}

    \begin{column}{0.42\textwidth}
      \visualplaceholder[5cm]{A multi-layer network solving XOR:
      input layer $\to$ hidden layer $\to$ output layer.
      The hidden layer creates a nonlinear boundary.
      Caption: ``Two layers CAN separate XOR.''}

      \vspace{0.3em}

      % Simple timeline of the drama arc
      \begin{tikzpicture}[font=\tiny, node distance=0.6cm]
        \node[fill=green!20, rounded corners, inner sep=2pt]
          (hype) {1958: Perceptron hype};
        \node[fill=red!20, rounded corners, inner sep=2pt, below=of hype]
          (crash) {1969: XOR crash};
        \node[fill=gray!20, rounded corners, inner sep=2pt, below=of crash]
          (winter) {1969--86: AI Winter};
        \node[fill=green!30, rounded corners, inner sep=2pt, below=of winter]
          (back) {1986: Backprop revival};
        \draw[-{Stealth[length=3pt]}, thick, green!60!black]
          (hype) -- (crash);
        \draw[-{Stealth[length=3pt]}, thick, red!60]
          (crash) -- (winter);
        \draw[-{Stealth[length=3pt]}, thick, green!60!black]
          (winter) -- (back);
      \end{tikzpicture}
    \end{column}
  \end{columns}

  \note{
    ``When backpropagation showed that multi-layer networks could learn,
    neural networks came roaring back.''

    The Rumelhart/Hinton/Williams 1986 paper was already covered in Slide~22
    (Section~2, Calculus). Cross-reference explicitly: ``Remember the chain
    rule? This is where it saved an entire field.''
  }
\end{frame}

% ----------------------------------------------------------
% Slide 59: Activation Functions
% ----------------------------------------------------------
\begin{frame}{Activation Functions --- The Spice That Makes It Work}

  \begin{columns}[T]
    \begin{column}{0.55\textwidth}
      \begin{insightbox}
        Without a nonlinear activation function, a multi-layer network
        collapses to a single layer. Activation functions give neural
        networks their power.
      \end{insightbox}

      \vspace{0.2em}
      {\small
      \textbf{Sigmoid} --- $\displaystyle\frac{1}{1+e^{-x}}$\par
      \begin{itemize}\setlength\itemsep{1pt}
        \item Smooth S-curve, used 1980s--2000s
        \item Related to the logistic function of
              \textbf{Pierre Fran\c{c}ois Verhulst} (1804--1849,
              Brussels), who modelled population growth (1838)
      \end{itemize}

      \textbf{ReLU} --- $\max(0, x)$\par
      \begin{itemize}\setlength\itemsep{1pt}
        \item Popularised by \textbf{Nair \& Hinton} (2010)
        \item Simpler, faster, avoids vanishing gradients
      \end{itemize}

      \textbf{GELU} --- used in GPT and BERT\par
      \begin{itemize}\setlength\itemsep{1pt}
        \item Introduced by \textbf{Hendrycks \& Gimpel} (2016)
        \item Smooth version of ReLU; used in modern transformers
      \end{itemize}
      }
    \end{column}

    \begin{column}{0.42\textwidth}
      \visualplaceholder[5cm]{Three activation function curves plotted
      side by side: Sigmoid (smooth S-curve, labelled ``1980s--2000s''),
      ReLU (hockey stick, labelled ``2010s''), GELU (smooth ReLU,
      labelled ``2016--present, used in GPT'').
      Caption: ``These simple curves make neural networks nonlinear
      --- and therefore powerful.''}
    \end{column}
  \end{columns}

  \note{
    ``Without activation functions, stacking 100 layers of matrix
    multiplication would be the same as one layer. The activation function
    adds the nonlinearity. It is a small thing that makes everything
    else work.''

    \verifytag{Verhulst logistic function dates --- 1838 or 1845 or both?}
    \verifytag{Nair/Hinton 2010 ICML paper as key popularisation; ReLU-like
    functions appeared earlier (Hahnloser et al.\ 2000).}
    \verifytag{Hendrycks/Gimpel 2016 GELU paper.}
  }
\end{frame}

% ----------------------------------------------------------
% Slide 60: Residual Connections (ResNet)
% ----------------------------------------------------------
\begin{frame}{Residual Connections --- The Trick That Enabled Depth}

  \begin{columns}[T]
    \begin{column}{0.52\textwidth}
      \begin{personbox}[{Kaiming He et al.\ (b.\,c.\,1984)}]{aiViolet}
        \textbf{Origin:} China / USA\\[3pt]
        \textbf{Key paper (2015):} \emph{``Deep Residual Learning for
        Image Recognition''} (ResNet).\\[3pt]
        Instead of each layer computing a completely new output,
        it computes a \alert{change} to the input:
        \[ \text{output} = \text{input} + \text{layer}(\text{input}) \]
        This simple idea allowed networks to grow from
        $\sim$20~layers to 152~layers --- and beyond.
      \end{personbox}

      \vspace{0.3em}
      \begin{insightbox}
        The ``\textbf{Add}'' in the transformer's ``Add \& Norm'' block
        IS He's residual connection from 2015.
        Without it, transformers could not be deep.
      \end{insightbox}
    \end{column}

    \begin{column}{0.45\textwidth}
      \visualplaceholder[5.5cm]{Side-by-side comparison: Left --- standard
      deep network with signal fading layer by layer (fading colour).
      Right --- ResNet with skip connections; an arrow skips over each block,
      carrying the original signal forward.
      Below: residual formula as diagram: input $\to$ [layer] $\to$
      ADD input back $\to$ output.}
    \end{column}
  \end{columns}

  \note{
    ``He's insight was deceptively simple: instead of asking each layer to
    compute the full answer, ask it to compute the DIFFERENCE from the
    input. If a layer has nothing useful to add, it can output zero and
    the input passes through unchanged.''

    LLM connection: ``Every transformer block has a residual connection.''
  }
\end{frame}

% ----------------------------------------------------------
% Slide 61: Layer Normalization and Dropout
% ----------------------------------------------------------
\begin{frame}{Layer Normalization \& Dropout --- Stabilising Deep Networks}

  \begin{columns}[T]
    \begin{column}{0.52\textwidth}
      \begin{personbox}[Ba/Kiros/Hinton -- Layer Norm (2016)]{aiViolet}
        Adjusts the outputs of each layer so they have a consistent
        mean and variance. Stabilises training and speeds convergence.
      \end{personbox}

      \vspace{0.3em}

      \begin{personbox}[Srivastava/Hinton et al. -- Dropout (2014)]{aiViolet}
        Randomly deactivates neurons during training, forcing the network
        to learn robust features rather than memorising.\\[3pt]
        \emph{Like studying for an exam by covering random notes --- you
        learn the material, not the page layout.}
      \end{personbox}

      \vspace{0.3em}
      {\small\textbf{Layer Norm recipe:}
        (1)~Compute the average of all neuron outputs in a layer.
        (2)~Compute how spread out they are.
        (3)~Subtract the average and divide by the spread.
        Now they are centred at zero with consistent scale.}
    \end{column}

    \begin{column}{0.45\textwidth}
      \visualplaceholder[5.5cm]{Top: Layer normalisation --- a row of neuron
      activations before (wildly varying) and after normalisation (centred,
      consistent). Caption: ``Layer Norm = rescale so nothing explodes.''
      Bottom: Dropout --- a network with some neurons crossed out.
      Caption: ``Dropout = randomly disable neurons during training.''}

      \vspace{0.3em}
      \begin{tcolorbox}[colback=aiViolet!5,colframe=aiViolet!60,
        boxrule=0.5pt,rounded corners,left=3pt,right=3pt,top=2pt,bottom=2pt]
        \small In the transformer's ``Add \& Norm'':\\
        \textbf{Add} = He's residual connection (Slide~60)\\
        \textbf{Norm} = Ba's layer normalisation
      \end{tcolorbox}
    \end{column}
  \end{columns}

  \note{
    ``Layer normalisation is the `Norm' in the transformer's `Add \& Norm.'
    Dropout is used during training. Both are invisible to the user but
    essential for making deep networks work.''

    \verifytag{Ba et al.\ 2016 --- arXiv preprint, not initially peer-reviewed.}
    \verifytag{Srivastava et al.\ 2014 JMLR as definitive paper; Hinton
    proposed the idea earlier in 2012.}
  }
\end{frame}

% ----------------------------------------------------------
% Slide 62: Convergence Diagram Update -- The Neuron & Depth
% ----------------------------------------------------------
\begin{frame}{Convergence Diagram --- The Neuron \& Depth Lights Up}

  \visualplaceholder[10cm]{Convergence diagram: ``The Neuron \& Depth'' node
  now lit in yellow/gold. 7 of 10 nodes active.
  Label: ``Activation + normalisation + residuals = deep networks.''}

  \vspace{0.3em}
  \begin{columns}[T]
    \begin{column}{0.48\textwidth}
      \begin{tcolorbox}[colback=aiViolet!5,colframe=aiViolet,
        fonttitle=\bfseries\small,title={Section 7 Summary},
        rounded corners,boxrule=0.8pt,left=3pt,right=3pt,top=3pt,bottom=3pt]
        \scriptsize
        \begin{tabular}{@{}ll@{}}
          McCulloch-Pitts (1943)    & First math neuron\\
          Rosenblatt (1958)         & Perceptron learns\\
          Minsky/Papert (1969)      & XOR, AI winter\\
          Rumelhart/Hinton/W.\ (1986) & Multi-layer revival\\
          Verhulst (1838) $\to$ Sigmoid & Population $\to$ activation\\
          Nair/Hinton (2010)        & ReLU\\
          Srivastava et al.\ (2014)  & Dropout\\
          He et al.\ (2015)         & Residual connections\\
          Ba et al.\ (2016)         & Layer normalisation\\
          Hendrycks/Gimpel (2016)   & GELU
        \end{tabular}
      \end{tcolorbox}
    \end{column}

    \begin{column}{0.48\textwidth}
      \begin{tcolorbox}[colback=accentOrange!8,colframe=accentOrange,
        fonttitle=\bfseries\small,
        title={\faArrowRight\ Connection to the Transformer},
        rounded corners,boxrule=0.8pt,left=3pt,right=3pt,top=3pt,bottom=3pt]
        \scriptsize
        The artificial neuron, with its activation function, is the basic
        unit. Residual connections, layer normalisation, and dropout allow
        transformers to be DEEP. There are billions of neurons in GPT,
        stacked 96+ layers.
      \end{tcolorbox}
    \end{column}
  \end{columns}

  \note{
    Transition: ``We can build neurons, stack them into deep layers, and
    train them. But language is SEQUENTIAL --- word order matters.
    How do you process a sequence? That took another set of inventions.''
  }
\end{frame}


% ============================================================
% SECTION 8: SEQUENCE MODELING (Slides 63--68)
% "The Path to Language"
% Timeline: 1906--2014
% ============================================================
\section{Sequence Modeling}

% ----------------------------------------------------------
% Slide 63: Section Divider
% ----------------------------------------------------------
{
\setbeamercolor{background canvas}{bg=aiViolet}
\begin{frame}[plain]
  \centering
  \vspace{1.8cm}
  \textcolor{white}{\Huge\bfseries Sequence Modeling}\\[0.5em]
  \textcolor{white}{\Large The Path to Language}\\[1em]
  \textcolor{white!80}{\large 1906--2014}\\[1.5em]
  \textcolor{white!60}{\normalsize Markov \(\cdot\) Baum \(\cdot\) Elman
    \(\cdot\) Hochreiter \& Schmidhuber \(\cdot\) Bahdanau}\\[1em]
  \textcolor{white!50}{\small Language is a sequence.
  Modeling sequences required 100 years of invention.}

  \note{
    Convergence diagram: Section 8 node pulsing.

    ``Language is sequential --- word order matters. `Dog bites man' is
    very different from `Man bites dog.' How do you process a sequence?''
  }
\end{frame}
}

% ----------------------------------------------------------
% Slide 64: Markov Chains
% ----------------------------------------------------------
\begin{frame}{Andrey Markov --- The Simplest Language Model}

  \begin{columns}[T]
    \begin{column}{0.52\textwidth}
      \begin{personbox}[Andrey Andreyevich Markov (1856--1922)]{aiViolet}
        \textbf{Origin:} Ryazan, Russia\\[3pt]
        Introduced \textbf{Markov chains} (1906) by analysing the
        alternation of vowels and consonants in Pushkin's poem
        \emph{Eugene Onegin}.\\[3pt]
        He was \emph{not} studying language --- he was disproving a
        claim that the law of large numbers required independence.
        But his method became the foundation of all sequence modelling.
      \end{personbox}

      \vspace{0.3em}
      \begin{insightbox}
        A Markov chain predicts the next state based \emph{only} on the
        current state, ignoring all history. Your phone's autocomplete
        is a descendant of this idea.
      \end{insightbox}
    \end{column}

    \begin{column}{0.45\textwidth}
      \visualplaceholder[5.5cm]{Top: State diagram --- states are words
      (``the'', ``cat'', ``sat'', ``on''), arrows labelled with
      transition probabilities.
      Bottom: Text generated by a Markov chain vs.\ text from an LLM.
      Markov text is grammatical locally but incoherent globally.}
    \end{column}
  \end{columns}

  \note{
    ``Markov analysed poetry to prove a point about probability theory.
    His chain model became the ancestor of every language model, including
    the one in your phone's autocomplete.''

    \verifytag{``Eugene Onegin'' analysis dated to 1913 in some sources,
    while the mathematical framework was published 1906. Clarify the
    distinction.}
  }
\end{frame}

% ----------------------------------------------------------
% Slide 65: Hidden Markov Models
% ----------------------------------------------------------
\begin{frame}{Hidden Markov Models --- Adding Depth}

  \begin{columns}[T]
    \begin{column}{0.52\textwidth}
      \begin{personbox}[Leonard Baum et al.\ (1960s--1972)]{aiViolet}
        \textbf{Origin:} USA\\[3pt]
        Developed the mathematical framework for \textbf{Hidden Markov
        Models} (HMMs) and the \textbf{Baum--Welch algorithm} for
        training them.\\[3pt]
        HMMs add a hidden state layer: the observed output depends on
        an unobserved internal state that must be inferred.
      \end{personbox}

      \vspace{0.3em}
      \begin{insightbox}
        Speech recognition used HMMs for 30+ years before deep learning
        replaced them. The idea of ``hidden internal states'' producing
        observable outputs is conceptually similar to the hidden layers
        of a neural network.
      \end{insightbox}
    \end{column}

    \begin{column}{0.45\textwidth}
      \visualplaceholder[5.5cm]{HMM diagram: hidden states on top (connected
      by transition arrows), observed outputs below (connected to hidden
      states by emission arrows).
      Caption: ``The hidden states cannot be seen directly --- they must
      be inferred from observations.''}
    \end{column}
  \end{columns}

  \note{
    ``HMMs dominated speech recognition from the 1970s through the 2010s.
    Siri's original speech engine used HMMs.''

    \verifytag{Baum--Welch algorithm dates: key papers 1966, 1970, 1972.
    Baum's full name and dates are surprisingly hard to pin down.}
  }
\end{frame}

% ----------------------------------------------------------
% Slide 66: Recurrent Neural Networks
% ----------------------------------------------------------
\begin{frame}{Recurrent Neural Networks --- Memory in a Loop}

  \begin{columns}[T]
    \begin{column}{0.52\textwidth}
      \begin{personbox}[Michael Jordan -- Jordan Network (1986)]{aiViolet}
        \textbf{Origin:} USA\\[3pt]
        An early form of recurrent neural network.\\
        \emph{Yes, the field has its own Michael Jordan.}
      \end{personbox}

      \vspace{0.2em}

      \begin{personbox}[Jeffrey Elman (1948--2018) --- Elman Network (1990)]{aiViolet}
        \textbf{Origin:} USA\\[3pt]
        The hidden state at time $t$ is fed back as input at time $t{+}1$.
        This is the standard form of RNN.
      \end{personbox}

      \vspace{0.2em}
      \begin{tcolorbox}[colback=red!5,colframe=red!60,
        boxrule=0.5pt,rounded corners,left=3pt,right=3pt,top=2pt,bottom=2pt]
        \small\textbf{Fatal flaw:} The hidden state gets overwritten
        at each step. After 50~words, the memory of word~1 has been
        washed away. This is the \alert{vanishing gradient problem}.
      \end{tcolorbox}
    \end{column}

    \begin{column}{0.45\textwidth}
      \visualplaceholder[5.5cm]{An RNN unrolled across time steps: the same
      network processing word 1, then word 2, then word 3, with the
      hidden state carried forward like a ``memory'' arrow.
      Caption: ``The hidden state is the RNN's memory. But it is a
      very leaky memory.''}
    \end{column}
  \end{columns}

  \note{
    ``RNNs have a fatal flaw: the hidden state gets overwritten at each step.
    By the time you have read 50 words, the memory of word 1 has been
    washed away. This is the vanishing gradient problem.''

    \verifytag{Jordan 1986 --- sometimes cited as 1986 technical report.}
  }
\end{frame}

% ----------------------------------------------------------
% Slide 67: LSTM [EXTENDED]
% ----------------------------------------------------------
\begin{frame}{LSTM --- Long Short-Term Memory
  \hfill\small\textcolor{gray}{[EXTENDED]}}

  \begin{columns}[T]
    \begin{column}{0.52\textwidth}
      \begin{personbox}[Hochreiter \& Schmidhuber]{aiViolet}
        \textbf{Sepp Hochreiter:} Austria\\
        \textbf{J\"urgen Schmidhuber:} Munich, Germany\\[3pt]
        Published \emph{``Long Short-Term Memory''} (1997).\\[3pt]
        Hochreiter's diploma thesis (1991, supervised by Schmidhuber)
        identified the vanishing gradient problem. The 1997 LSTM paper
        provided the solution: explicit \textbf{memory cells} with
        \textbf{gates} that control what to remember and what to forget.
      \end{personbox}

      \vspace{0.3em}
      \begin{insightbox}
        LSTM dominated language tasks from the late 1990s through 2017.
        Google Translate used LSTMs. But LSTMs process words
        \emph{one at a time} --- you cannot parallelise them.
        This bottleneck led to \alert{attention}.
      \end{insightbox}
    \end{column}

    \begin{column}{0.45\textwidth}
      \visualplaceholder[5.5cm]{Simplified LSTM cell: input gate, forget
      gate, output gate, cell state. Caption: ``The forget gate decides
      what to remember. The input gate decides what is new. The output
      gate decides what to reveal.''}
    \end{column}
  \end{columns}

  \note{
    ``LSTM solved the memory problem. But it introduced a new bottleneck:
    sequential processing. You cannot parallelise an LSTM --- each word
    must wait for the previous one. For long documents, this is painfully
    slow. The key idea that changed everything: attention.''

    This slide is EXTENDED --- can cut if short on time. Mention LSTM
    en route to attention if skipping.

    \verifytag{Hochreiter 1991 diploma thesis as vanishing gradient analysis.}
  }
\end{frame}

% ----------------------------------------------------------
% Slide 68: The Attention Mechanism -- KEY BREAKTHROUGH
% ----------------------------------------------------------
\begin{frame}{The Attention Mechanism --- The Bridge to Transformers}

  \begin{columns}[T]
    \begin{column}{0.52\textwidth}
      \begin{personbox}[Bahdanau/Cho/Bengio (2014)]{aiViolet}
        \textbf{Dzmitry Bahdanau:} Belarus/Canada\\
        \textbf{Kyunghyun Cho:} South Korea/Canada\\
        \textbf{Yoshua Bengio:} France/Canada\\[3pt]
        Published \emph{``Neural Machine Translation by Jointly Learning
        to Align and Translate''} (2014 arXiv; 2015 ICLR).\\[3pt]
        \textbf{Key insight:} Instead of compressing the entire input
        into a fixed-size vector, let the model \alert{look back}
        at all previous positions and decide which ones are relevant.
      \end{personbox}

      \vspace{0.3em}
      \begin{tcolorbox}[colback=aiViolet!8,colframe=aiViolet,
        fonttitle=\bfseries,title={The Breakthrough},
        rounded corners,boxrule=1pt,left=4pt,right=4pt,top=3pt,bottom=3pt]
        \small \textbf{Attention} = the ability to look back and focus.\\[3pt]
        Without this 2014 paper, there would be no GPT,
        no Claude, no ChatGPT.
      \end{tcolorbox}
    \end{column}

    \begin{column}{0.45\textwidth}
      \visualplaceholder[5.5cm]{Top: Attention diagram --- translating
      ``Le chat est sur le tapis.'' When generating ``cat,'' the model
      attends strongly to ``chat.'' When generating ``mat,'' it attends
      to ``tapis.'' Colour-coded attention weights shown as lines of
      varying thickness.
      Bottom: Comparison --- RNN (compresses everything into one vector)
      vs.\ Attention (can look back at everything).}
    \end{column}
  \end{columns}

  \note{
    ``Bahdanau's insight was simple but transformative: why compress
    everything into one vector when you can look back at the whole input?
    This idea --- attention --- became the core of the transformer.''

    Give this slide weight. It is the key intellectual bridge.

    Convergence diagram update: ``Sequence Modeling'' node lit in teal.
    Label: ``Attention = context.''

    Transition: ``We now have ALL the building blocks. It is time to
    put them together.''
  }
\end{frame}


% ============================================================
% SECTION 9: THE TRANSFORMER (Slides 69--74)
% "Where Everything Converges"
% Timeline: 2017
% ============================================================
\section{The Transformer}

% ----------------------------------------------------------
% Slide 69: Section Divider
% ----------------------------------------------------------
{
\setbeamercolor{background canvas}{bg=aiViolet}
\begin{frame}[plain]
  \centering
  \vspace{1.6cm}
  \textcolor{white}{\Huge\bfseries The Transformer}\\[0.5em]
  \textcolor{white}{\Large Where Everything Converges}\\[1em]
  \textcolor{white!80}{\large 2017}\\[1.5em]
  \textcolor{white!60}{\normalsize Vaswani \(\cdot\) Shazeer \(\cdot\)
    Parmar \(\cdot\) Uszkoreit \(\cdot\) Jones}\\[0.2em]
  \textcolor{white!60}{\normalsize Gomez \(\cdot\) Kaiser \(\cdot\)
    Polosukhin}\\[1em]
  \textcolor{white!50}{\small Eight branches of mathematics meet
  in one architecture.}\\[0.3em]
  \textcolor{white!40}{\footnotesize\itshape ``Attention Is All You Need.''}

  \note{
    Convergence diagram: 8 of 10 nodes lit. Section 9 pulsing.

    ``2017: Eight branches of mathematics, developed over more than two
    millennia, all converge in one architecture.''
  }
\end{frame}
}

% ----------------------------------------------------------
% Slide 70: "Attention Is All You Need"
% ----------------------------------------------------------
\begin{frame}{``Attention Is All You Need'' --- The Paper That Changed AI}

  \begin{columns}[T]
    \begin{column}{0.55\textwidth}
      {\small In 2017, a team at Google published a paper proposing an
      architecture that used \textbf{only attention} --- no recurrence,
      no convolution. It was simpler, faster, and better.}

      \vspace{0.4em}

      \begin{personbox}[The Eight Authors (2017 NeurIPS)]{aiViolet}
        \textbf{Ashish Vaswani}, \textbf{Noam Shazeer},
        \textbf{Niki Parmar}, \textbf{Jakob Uszkoreit},
        \textbf{Llion Jones}, \textbf{Aidan Gomez},
        \textbf{{\L}ukasz Kaiser}, \textbf{Illia Polosukhin}\\[3pt]
        Most in their 20s and 30s. Several have since founded major AI
        companies (Cohere, Character.AI).
      \end{personbox}

      \vspace{0.3em}
      \begin{insightbox}
        Eight researchers. One paper. The architecture that powers
        \emph{every} major LLM.
      \end{insightbox}
    \end{column}

    \begin{column}{0.42\textwidth}
      \visualplaceholder[5cm]{Top: The first page of ``Attention Is All You
      Need'' with title highlighted.
      Bottom: Stylised portraits or photos of the eight authors.
      Caption: ``Eight researchers. One paper. One architecture.''}
    \end{column}
  \end{columns}

  \note{
    ``The title is bold: `Attention Is All You Need.' They were right.
    This single architecture --- the transformer --- is the foundation of
    GPT, Claude, Gemini, and virtually every LLM in existence.''

    \verifytag{Gomez founded Cohere; Jones and Shazeer co-founded
    Character.AI. Confirm details.}
  }
\end{frame}

% ----------------------------------------------------------
% Slide 71: Self-Attention -- The Core Mechanism
% ----------------------------------------------------------
\begin{frame}{Self-Attention --- The Core Mechanism}

  {\small Each token generates three vectors: a \textbf{Query}
  (``what am I looking for?''), a \textbf{Key} (``what do I contain?''),
  and a \textbf{Value} (``what information do I carry?'').}

  \vspace{0.3em}

  % Step-by-step with back-references
  \begin{tcolorbox}[colback=aiViolet!5,colframe=aiViolet,
    fonttitle=\bfseries,title={Self-Attention Step by Step},
    rounded corners,boxrule=0.8pt,left=4pt,right=4pt,top=4pt,bottom=4pt]
    \small
    \begin{enumerate}\setlength\itemsep{2pt}
      \item Compute attention scores: \textbf{dot product} of each Query
            with each Key \hfill\textcolor{gray}{\scriptsize(Gibbs, Slide~10)}
      \item Scale scores: divide by $\sqrt{d}$ to keep numbers
            manageable \hfill\textcolor{gray}{\scriptsize(cosine-like scaling, Slide~11)}
      \item Normalise with \textbf{softmax}: scores become
            probabilities \hfill\textcolor{gray}{\scriptsize(Boltzmann, Slide~33)}
      \item Output = weighted sum of Values via
            \textbf{matrix multiplication}
            \hfill\textcolor{gray}{\scriptsize(Cayley, Slide~9)}
    \end{enumerate}
  \end{tcolorbox}

  \vspace{0.2em}

  % Formal inset
  \begin{columns}[T]
    \begin{column}{0.48\textwidth}
      \begin{tcolorbox}[colback=gray!5,colframe=gray!40,
        boxrule=0.5pt,rounded corners,left=3pt,right=3pt,top=2pt,bottom=2pt]
        {\scriptsize\textbf{Formal notation (small inset):}
        \[
          \text{Attention}(Q,K,V)
            = \text{softmax}\!\left(\frac{Q K^{\!\top}}{\sqrt{d_k}}\right) V
        \]
        $K^{\!\top}$ means ``flip the matrix on its side'' (transpose).}
      \end{tcolorbox}
    \end{column}
    \begin{column}{0.48\textwidth}
      {\small\textbf{Concrete example:} In ``The cat sat on the mat because
      \emph{it} was tired,'' the word ``it'' attends most strongly to
      ``cat.'' Self-attention learns that ``it'' refers to ``cat.''}
    \end{column}
  \end{columns}

  \note{
    ``Self-attention is the convergence point. Linear algebra, probability,
    and information theory, all in one equation.''

    Walk through the concrete ``it'' $\to$ ``cat'' example. This is how
    the model resolves references across a sentence.
  }
\end{frame}

% ----------------------------------------------------------
% Slide 72: Feed-Forward Block and Multi-Head Attention
% ----------------------------------------------------------
\begin{frame}{Multi-Head Attention \& Feed-Forward Networks}

  \begin{columns}[T]
    \begin{column}{0.52\textwidth}
      \textbf{Multi-Head Attention}\\[3pt]
      {\small Instead of one set of Q/K/V, the transformer uses multiple
      ``heads'' (8, 12, or more) in parallel. Each head learns different
      relationships: syntax, semantics, position, etc.
      Results are concatenated and projected back.}

      \vspace{0.5em}
      \textbf{Feed-Forward Network (FFN)}\\[3pt]
      {\small Two matrix multiplications with a GELU activation in between.
      This is the universal approximation theorem
      (Cybenko, Slide~52) applied within each transformer layer ---
      the FFN provides the raw computational power.}

      \vspace{0.5em}
      \begin{insightbox}
        \small One transformer layer = Attention + FFN + two rounds of
        ``Add \& Norm.'' Stack 96 of these and you have GPT-4.
      \end{insightbox}
    \end{column}

    \begin{column}{0.45\textwidth}
      \visualplaceholder[5.5cm]{A single transformer block, annotated:
      Multi-Head Attention sub-block $\to$ Add \& Norm
      (residual + layer norm, Slides 60--61) $\to$
      Feed-Forward Network sub-block $\to$ Add \& Norm.
      Every component labelled with its origin section.}
    \end{column}
  \end{columns}

  \note{
    ``One transformer layer = attention + feed-forward + two rounds of
    `add and normalise.' Stack 96 of these layers and you have GPT-4.
    Every component was taught in an earlier section of this lecture.''
  }
\end{frame}

% ----------------------------------------------------------
% Slide 73: THE CLIMAX -- Full Annotated Transformer Diagram
% ----------------------------------------------------------
\begin{frame}{The Full Transformer --- Annotated with 2,000+ Years of Math}

  % Large placeholder -- THIS IS THE LECTURE'S CLIMAX
  \visualplaceholder[13cm]{%
    \textbf{CLIMAX DIAGRAM --- Full decoder-only transformer architecture}\\[4pt]
    Annotated components (back-references to earlier slides):\\
    \textbullet\ \textbf{Input Embedding:} Words $\to$ Vectors
      (Grassmann 1844, Mikolov 2013, Slide~14)\\
    \textbullet\ \textbf{Positional Encoding:} Sine/cosine
      (Fourier 1807, Slide~50)\\
    \textbullet\ \textbf{Multi-Head Self-Attention:} Dot product + softmax
      + matrix mult.\ (Gibbs 1880s Slide~10, Boltzmann 1870s Slide~33,
      Cayley 1858 Slide~9, Bahdanau 2014 Slide~68)\\
    \textbullet\ \textbf{Add \& Norm:} Residual + Layer Norm
      (He 2015 Slide~60, Ba 2016 Slide~61)\\
    \textbullet\ \textbf{FFN:} Matrix mult.\ + GELU
      (Cybenko 1989 Slide~52, Hendrycks 2016 Slide~59)\\
    \textbullet\ \textbf{Dropout:} (Srivastava/Hinton 2014 Slide~61)\\
    \textbullet\ \textbf{Output Softmax:} Probability over vocabulary
      (Fisher 1922, Boltzmann/Luce Slide~33)\\
    \textbullet\ \textbf{Training:} Cross-entropy + backprop + SGD
      (Shannon 1948, Cauchy 1847, Linnainmaa 1970)\\
    \textbullet\ \textbf{Hardware:} GPU arrays
      (NVIDIA CUDA 2007, AlexNet 2012 Slide~47)%
  }

  \vspace{0.3em}
  \begin{tcolorbox}[colback=accentOrange!10,colframe=accentOrange,
    fonttitle=\bfseries,
    title={\faUsers\ Interactive Moment: ``Count the Building Blocks''},
    rounded corners,boxrule=1pt,left=4pt,right=4pt,top=3pt,bottom=3pt]
    \small Students: look at this diagram and try to identify which of
    the 10~sections each component comes from. How many can you find?
    \hfill\emph{(3~min)}
  \end{tcolorbox}

  \note{
    THIS IS THE CENTRAL THESIS OF THE ENTIRE LECTURE, crystallised in
    one diagram. Let this land. \textbf{Pause.}

    ``Look at this diagram. Every single piece has a history. The sine
    functions come from Fourier in 1807. The matrix multiplications come
    from Cayley in 1858. The softmax comes from Boltzmann in the 1870s.
    The training algorithm uses Cauchy's gradient descent from 1847 and
    Linnainmaa's automatic differentiation from 1970. The residual
    connections come from He in 2015. The layer normalisation from Ba in
    2016. When you talk to an AI, you are talking to 2,000 years of
    mathematics, all running at once.''

    Interactive moment: give students 3 minutes to ``count the building
    blocks'' --- identify which of the 10 sections each component comes
    from.
  }
\end{frame}

% ----------------------------------------------------------
% Slide 74: Positional Encoding -- Fourier Meets Transformers
% ----------------------------------------------------------
\begin{frame}{Positional Encoding --- Fourier Meets Transformers}

  \begin{columns}[T]
    \begin{column}{0.52\textwidth}
      {\small Transformers have no built-in notion of word order (unlike
      RNNs). To fix this, they add \textbf{positional encodings} ---
      sine and cosine functions at different frequencies.
      This is Fourier's idea, applied directly.}

      \vspace{0.4em}

      \textbf{Recipe (CONCEPTUAL):}\\
      {\small For each position in the sentence and each dimension in
      the vector, compute a sine or cosine wave with a specific
      frequency. Different dimensions use different frequencies ---
      just like Fourier's decomposition. Each position gets a unique
      ``fingerprint'' of wave values.}

      \vspace{0.3em}

      % Formal inset
      \begin{tcolorbox}[colback=gray!5,colframe=gray!40,
        boxrule=0.5pt,rounded corners,left=3pt,right=3pt,top=2pt,bottom=2pt]
        {\scriptsize\textbf{Formal:}
        \begin{align*}
          PE(\text{pos}, 2i) &= \sin\!\bigl(\text{pos}\, /\, 10000^{2i/d}\bigr)\\
          PE(\text{pos}, 2i{+}1) &= \cos\!\bigl(\text{pos}\, /\, 10000^{2i/d}\bigr)
        \end{align*}
        The $10000^{2i/d}$ creates a spectrum of frequencies from
        very slow to very fast.}
      \end{tcolorbox}
    \end{column}

    \begin{column}{0.45\textwidth}
      \visualplaceholder[5.5cm]{Heatmap of positional encodings:
      rows = positions (1st word, 2nd word, \ldots, 100th word),
      columns = dimensions. Clear sinusoidal wave patterns visible.
      Caption: ``Fourier's sine and cosine waves, used to tell
      the transformer which word is which.''}

      \vspace{0.3em}
      \begin{insightbox}
        Fourier decomposed heat into waves.
        Vaswani decomposed word position into waves.
        Same math, 210 years apart.
      \end{insightbox}
    \end{column}
  \end{columns}

  \note{
    ``Fourier decomposed heat into waves. Vaswani decomposed word
    position into waves. The same math, 210 years apart, for completely
    different purposes.''

    Convergence diagram update: ALL 10 nodes now lit (Sections 9 and 10).
    Central ``?'' becomes ``TRANSFORMER.'' The full convergence diagram
    is complete.
  }
\end{frame}


% ============================================================
% SECTION 10: SCALING TO LLMs + CLOSING (Slides 75--78)
% "The Final Leap"
% Timeline: 2018--2026
% ============================================================
\section{Scaling to LLMs}

% ----------------------------------------------------------
% Slide 75: The Scaling Leap
% ----------------------------------------------------------
\begin{frame}{The Scaling Leap --- From Transformer to LLM}

  \begin{columns}[T]
    \begin{column}{0.52\textwidth}
      {\small The transformer architecture, scaled up with more parameters,
      more data, and more compute, produces Large Language Models with
      emergent abilities.}

      \vspace{0.3em}

      \begin{tcolorbox}[colback=aiViolet!5,colframe=aiViolet,
        fonttitle=\bfseries\small,title={The GPT Series},
        rounded corners,boxrule=0.8pt,left=3pt,right=3pt,top=3pt,bottom=3pt]
        \scriptsize
        \begin{tabular}{@{}lrl@{}}
          \textbf{GPT-1} (2018) & 117M params & Pre-train + fine-tune\\
          \textbf{GPT-2} (2019) & 1.5B params & Convincing text gen.\\
          \textbf{GPT-3} (2020) & 175B params & Few-shot learning\\
          \textbf{GPT-4} (2023) & undisclosed & Multimodal\\
          \textbf{Claude} (2023--26) & --- & Constitutional AI
        \end{tabular}
      \end{tcolorbox}

      \vspace{0.3em}

      \begin{personbox}[Jared Kaplan et al.\ (2020) --- Scaling Laws]{aiViolet}
        \scriptsize LLM performance scales as a smooth \textbf{power law}
        with compute, data, and parameters. Predictable improvement.
      \end{personbox}
    \end{column}

    \begin{column}{0.45\textwidth}
      \visualplaceholder[5.5cm]{Top: Log-log plot of parameter count vs.\
      year: GPT-1 (2018) $\to$ GPT-2 (2019) $\to$ GPT-3 (2020) $\to$
      PaLM (2022) $\to$ GPT-4 (2023) $\to$ frontier models (2025--26).
      Exponential growth visible.
      Bottom: Timeline of key milestones: RLHF (2017), InstructGPT (2022),
      ChatGPT (Nov 2022), GPT-4 (Mar 2023), Claude~3 (2024), open-source
      explosion (Llama, Mistral).}
    \end{column}
  \end{columns}

  \note{
    ``The transformer is the same architecture from 2017. What changed is
    scale. More parameters, more training data, more compute. And somewhere
    along the way, these models started doing things nobody programmed them
    to do --- reasoning, writing code, translating languages they were
    barely trained on.''

    \verifytag{Kaplan et al.\ 2020 scaling laws paper.}
    \verifytag{ChatGPT launch: November 30, 2022.}
    \verifytag{Christiano et al.\ 2017 for RLHF.}
  }
\end{frame}

% ----------------------------------------------------------
% Slide 76: Return to the Exploded Diagram
% ----------------------------------------------------------
\begin{frame}{Return to the Exploded Diagram --- Now You Understand Every Piece}

  \visualplaceholder[13cm]{%
    \textbf{The SAME exploded transformer diagram from Slide~3},
    now with every label including the mathematician's name, date,
    AND slide number:\\[3pt]
    \textbullet\ Matrix multiplication (Cayley, 1858 --- Slide~9)
    \textbullet\ Dot product (Gibbs, 1880s --- Slide~10)\\
    \textbullet\ Gradient descent (Cauchy, 1847 --- Slide~21)
    \textbullet\ Softmax (Boltzmann, 1870s --- Slide~33)\\
    \textbullet\ Cross-entropy (Shannon, 1948 --- Slide~37)
    \textbullet\ Backpropagation (Linnainmaa 1970, Rumelhart/Hinton/Williams
    1986 --- Slide~22)\\
    \textbullet\ Activation functions (Verhulst 1838, Nair/Hinton 2010
    --- Slide~59)\\
    \textbullet\ Residual connections (He et al.\ 2015 --- Slide~60)
    \textbullet\ Layer norm (Ba et al.\ 2016 --- Slide~61)\\
    \textbullet\ Self-attention (Bahdanau 2014, Vaswani 2017
    --- Slides~68, 71)\\
    \textbullet\ Positional encoding (Fourier 1807 --- Slide~74)
    \textbullet\ GPU computing (NVIDIA CUDA 2007 --- Slide~47)%
  }

  \note{
    ``At the start of this lecture, this diagram was overwhelming.
    Now you can point to each piece and say who invented it, when,
    and why. That is the power of understanding the history.''

    Caption on slide: ``You now know the history of every piece of
    this machine.''
  }
\end{frame}

% ----------------------------------------------------------
% Slide 77: Convergence Diagram -- Final State
% ----------------------------------------------------------
\begin{frame}{Convergence Diagram --- Final State}

  \visualplaceholder[10cm]{The completed convergence diagram. All 10 nodes
  lit. Central ``TRANSFORMER'' node glowing. Each node label includes its
  cumulative date range:\\
  1.\ Linear Algebra (1809--2013) --- blue\\
  2.\ Calculus \& Optimisation (c.\,250 BCE--1986) --- green\\
  3.\ Probability \& Statistics (1654--1959) --- orange\\
  4.\ Information Theory (1865--2016) --- red\\
  5.\ Logic \& Computation (1854--2012) --- grey\\
  6.\ Function Approximation (1807--1991) --- purple\\
  7.\ The Neuron \& Depth (1943--2016) --- yellow\\
  8.\ Sequence Modeling (1906--2014) --- teal\\
  9.\ The Transformer (2017) --- multicolour\\
  10.\ Scaling to LLMs (2018--2026) --- gradient}

  \vspace{0.3em}
  \begin{center}
    {\large\bfseries\color{aiViolet}
    Ten streams of human invention, flowing into one machine.}
  \end{center}

  \note{
    Let the diagram speak. This is the visual summary of the entire lecture.
    All 10 nodes active. All arrows visible. 2,000+ years of mathematics
    converging.
  }
\end{frame}

% ----------------------------------------------------------
% Slide 78: Closing -- The Story Is Not Over
% ----------------------------------------------------------
\begin{frame}{The Story Is Not Over}

  \begin{tcolorbox}[colback=aiViolet!8,colframe=aiViolet,
    fonttitle=\bfseries\Large,
    title={\raisebox{-0.2em}{\faInfinity}\ 2,000+ Years of Mathematics.
    One Machine. Your Turn.},
    rounded corners,boxrule=1.5pt,
    left=6pt,right=6pt,top=6pt,bottom=6pt]

    \begin{itemize}\setlength\itemsep{6pt}
      \item Every component of an LLM was invented by someone, for a
            specific purpose, often centuries before AI existed.

      \item Grassmann imagined abstract vector spaces. Boltzmann studied
            gas molecules. Fourier solved the heat equation. Markov
            analysed poetry. Archimedes approximated circles.
            \textbf{None of them imagined AI.}

      \item The math you learn in school --- algebra, calculus,
            probability --- is the \emph{same math} running inside the
            AI you talk to every day.

      \item The next building block has not been invented yet.
            It might come from pure mathematics, from physics, from
            linguistics, from neuroscience ---
            \textbf{or from someone in this room.}
    \end{itemize}
  \end{tcolorbox}

  \vspace{0.3em}
  \visualplaceholder[10cm]{Reprise of the completed convergence diagram
  from Slide~77 / Exploded diagram from Slide~3.
  Below: ``2,000+ years of mathematics. One machine. Your turn.''}

  \note{
    ``When you talk to an AI, you are talking to the accumulated genius
    of Archimedes, Grassmann, Cayley, Newton, Leibniz, Cauchy, Boltzmann,
    Shannon, Turing, Mikolov, McCulloch, Pitts, Rosenblatt, Hochreiter,
    Schmidhuber, Bahdanau, Vaswani, and dozens of others. Two thousand
    years of human curiosity, all running at once, on hardware that
    Boole's algebra made possible. The next invention is yours.''

    Let the final line resonate. End.
  }
\end{frame}
